\section{Implementation Details}
\label{ch:implementation}

The experiments in this study are implemented in Python using PyTorch and the Hugging Face Transformers library. 
The pretrained weights for both NX-AI/xLSTM-7B and mistralai/Mistral-7B-v0.3 are obtained directly from the Hugging
Face model hub, ensuring a consistent and reproducible experimental setup. \\ \\
%
For benchmark evaluation, we use the LM Evaluation Harness <cite> framework as the primary tool for running standardized model assessments. 
Both this benchmark and the question-answering experiments require multiple inference operations, and for this reason we employ 
NVIDIA A100 GPUs with 40GB and 80GB of memory. The computation is executed in a cloud environment using the Modal infrastructure.\\ \\
%
A significant part of the study relies on accessing internal model representations. For all experiments involving hidden states, 
we employ PyTorch forward hooks to extract intermediate activations during inference. 
This mechanism allows us to capture layer-wise representations without modifying the original model architecture.\\ \\
%
When analyzing the cell states and input/forget gates, these values were not directly 
exposed by the model API. To address this limitation, we reconstruct them by intercepting the hidden states 
before and after each mLSTM block using hooks, and computing the corresponding values using the mathematical 
formulations provided in the xLSTM paper, as well as the reference implementation in the Hugging Face Transformers codebase. \\ \\
%
For the memory caching experiment, we build upon an existing xLSTM caching implementation provided in the model repository
,adapting it to support our experimental setup. This allows us to reuse previously computed memory states during question answering,
reducing redundant computation during inference.
We restrict decoding to a maximum of 25 new tokens and use a temperature of 1.0 
to ensure consistent sampling behavior across runs. The limit of 25 generated tokens is chosen because the SQuAD dataset 
primarily requires short, span-like answers, making longer generations unnecessary for the task. 
Inference time is measured by aggregating per-step execution intervals from the start to the end of generation. \\ \\ 
%
Finally, it is important to note that the xLSTM implementation benefits from optimized custom Triton kernels, 
which provide a significant speedup compared to standard PyTorch kernels. This optimization has a direct impact 
on overall runtime performance across all experiments.